% erdos634_v3.tex -- version 3.0 (2026-10-19).  Self-contained companion to the version 2.2 paper (erdos634_primes_and_N200.tex).
% Labels (Rule 0): \PL = PROVED-in-Lean (kernel-checked, standard axioms only, no sorry); \PH = PROVED-by-hand (written proof, audited, not formalized);
%   \CE = COMPUTED-exact (exact finite computation or exact verification of an explicit tiling, not a Lean proof); \CP = CLAIMED-by-Paliy (his public claim, not re-proved here);
%   \CJ = CONJECTURE.
% Figures: each \FIGSLOT{name} is a placeholder to be replaced by the figure of that name; the caption text is final.
\documentclass[11pt,reqno]{amsart}
\usepackage[T1]{fontenc}
\usepackage{amsmath,amssymb,amsthm}
\usepackage[margin=1.1in]{geometry}
\usepackage{array}
\usepackage{longtable}
\usepackage{booktabs}
\usepackage{graphicx}
\usepackage{hyperref}
\hypersetup{colorlinks=true,linkcolor=blue,citecolor=blue,urlcolor=blue}

\theoremstyle{plain}
\newtheorem{theorem}{Theorem}[section]
\newtheorem{lemma}[theorem]{Lemma}
\newtheorem{corollary}[theorem]{Corollary}
\newtheorem{proposition}[theorem]{Proposition}
\newtheorem{conjecture}[theorem]{Conjecture}
\theoremstyle{definition}
\newtheorem{definition}[theorem]{Definition}
\newtheorem{remark}[theorem]{Remark}
\newtheorem{example}[theorem]{Example}

\newcommand{\R}{\mathbb{R}}
\newcommand{\Z}{\mathbb{Z}}
\newcommand{\Q}{\mathbb{Q}}
\newcommand{\N}{\mathbb{N}}
\newcommand{\lab}[1]{\textup{[\textsc{#1}]}}
\newcommand{\PL}{\lab{proved in Lean}}
\newcommand{\PH}{\lab{proved by hand}}
\newcommand{\CE}{\lab{computed, exact}}
\newcommand{\CP}{\lab{claimed by Paliy}}
\newcommand{\CJ}{\lab{conjecture}}
\newcommand{\why}[1]{\par\smallskip\noindent\emph{Why it matters.} #1\par}
\DeclareUrlCommand{\lean}{\urlstyle{tt}}
\newcommand{\FIGSLOT}[1]{\begin{center}\fbox{\parbox{0.88\linewidth}{\centering\vspace{3.5em}\textsf{[figure slot: #1]}\vspace{3.5em}}}\end{center}}
\newcommand{\tile}{\mathsf{t}}
\newcommand{\kap}{\kappa}
\emergencystretch=2em

\title[Thresholds and a classification conjecture]{Congruent dissections of a triangle:\\ the six-piece template, scale thresholds, and a conjectured classification of the tile counts}
\author{Vico Bonfioli}
\date{Version 3.0, 19 October 2026}

\begin{document}

%=====================================================================
\begin{abstract}
A positive integer $N$ is a \emph{tile count} if some triangle can be cut into $N$ pairwise congruent triangles; Erd\H{o}s asked for the set of tile counts (Problem 634 of \texttt{erdosproblems.com}).
This problem is not solved here, and nothing in this paper bears on the prime case, which is treated in the earlier versions.
The paper is about the \emph{existence} side: which counts are realised, and at which scales.
Fix coprime integers $1\le e<f$ and the tile with sides $(ef,\,f^2-e^2,\,f^2)$; its angles $\alpha,\beta,\gamma$ satisfy $3\alpha+2\beta=\pi$.
D.~Paliy and M.~Beeson had shown that the triangles $W_M(e,f)$ and $T_M(e,f)$ built from this tile are tiled for every scale $M\ge f$, and the expectation was that nothing smaller works.
(1) \textbf{Theorem~W}: for $2\le e\le f-2$, $W_M(e,f)$ is tiled by $M^2(2f^2-e^2)$ tiles, and $T_M(e,f)$ by $M^2(3f^2-e^2)$ tiles, as soon as $Me\ge f$ and $(M-1)(f-e)\ge f$.
The proof is a layout of six pieces whose sizes are explicit linear functions of $M$ and three integers; it is formalised in Lean~4 for all $(e,f,M)$.
For $e<f/2$ the threshold $Me\ge f$ is sharp, so the scale $f$ is not the threshold in general; this answers the open case $W_4(3,5)$ of the earlier paper and refutes a size conjecture.
(2) A collar construction for the isosceles target with base angles $\alpha$ gives, in Lean, the counts $(f^2-1)(2f^2-1)g^2$ for all $g\ge f+1$ ($f\ge2$) and gives the tails $465\,g^2$ for $g\ge5$ and $966\,g^2$, $2485\,g^2$ for $g\ge2$; the up-set in $g$ was announced earlier by Paliy.
(3) Twenty-seven further tile counts are Lean theorems, by macro-piece layouts found by a template search and verified exactly; ninety-seven tilings are shipped as exact data.
(4) \textbf{Conjecture~C} states the set of tile counts as a union of explicit families; its existence half is a Lean theorem for the ranges proved so far, and its exclusion half is open.
For the non-squarefree $N\le10^6$ the census is now $199\,140$ excluded, $129\,108$ tile counts, $63\,826$ undecided; $770$ squarefree values stay open.
A list of corrections of earlier statements, and a list of what is not proved, are given.
\end{abstract}
\maketitle

%=====================================================================
\section{Introduction}\label{sec:intro}

\subsection{The problem and what this paper adds}
A triangle $\Omega$ (the \emph{target}) is \emph{tiled by $N$ copies of a triangle $\tile$} (the \emph{tile}) if $\Omega$ is the union of $N$ closed triangles with pairwise disjoint interiors, each congruent to $\tile$; mirror images are allowed and the pieces need not meet edge to edge.
$N$ is a \emph{tile count} if some triangle is tiled by $N$ copies of some triangle.
The squares, the numbers $x^2+y^2$, $3k^2$ and $6k^2$ are tile counts for classical reasons \cite{Soifer,SWW}.
Laczkovich classified the pairs (target, tile) up to the arithmetic of their angles \cite{Laczkovich1995,Laczkovich2012}; Beeson, with Ro\c{s}u, Zhang and Laczkovich, worked through the classes \cite{BeesonI,BeesonIII,BeesonIso,BeesonRosu,Zhang,Erdos633}.
Paliy published constructions and a proof of the prime case \cite{PaliyPrime,PaliyRepo}.
The earlier versions of this paper \cite{Release22,ReleasePaper} contain the exclusion side (the primes $p\equiv3\pmod4$, $p>3$, are not tile counts; every $N\le200$ is decided; a census of the non-squarefree $N\le10^6$).
Version~3.0 is about the other side, \emph{which counts are tile counts}, and it is organised around one family of tiles, the one that carries every open question of the census.

\medskip\noindent\textbf{Standing of the statements.}
Every statement carries one of five labels.
\PL{} The statement is a theorem of the accompanying Lean~4 library, checked by the kernel with the standard axioms only (\texttt{propext}, \texttt{Classical.choice}, \texttt{Quot.sound}), no \texttt{sorry}; the Lean statement has been compared with the mathematical one \emph{within this work only} (a reading by a person outside the project has not been made).
\PH{} A written proof, audited inside the project, not formalised.
\CE{} An exact finite computation, typically an explicit tiling whose pieces are listed in exact arithmetic and verified by an independent exact checker; this is evidence of existence, and for an explicit tiling it is a proof, but it is not a Lean proof.
\CP{} A public claim of D.~Paliy that we have not re-proved (in some cases we have checked an example).
\CJ{} Unproved.
The labels are repeated in the tables; where a result has two labels (say, a hand proof and a Lean proof) both are given.

\subsection{Results at a glance}
\begin{center}
\small
\begin{tabular}{@{}p{3.1cm}p{6.2cm}p{3.5cm}@{}}
\toprule
result & statement in one line & label and place\\
\midrule
Theorem W (\S\ref{sec:thmW}) & $W_M(e,f)$ and $T_M(e,f)$ tile once $Me\ge f$ and $(M-1)(f-e)\ge f$ ($2\le e\le f-2$) & \PL{} \PH{}; \S\ref{sec:thmW}, App.~\ref{app:W}\\
Sharpness (Cor.~\ref{cor:sharp}) & for $2\le e<f/2$: $W_M(e,f)$ tiles iff $Me\ge f$ & \PL{}\\
Collar family (\S\ref{sec:collar}) & $(f^2-1)(2f^2-1)g^2$, $g\ge f+1$; tails $465g^2$ $(g\ge5)$, $966g^2$ and $2485g^2$ $(g\ge2)$ & \PL{}; collar step with arbitrary height \PH{}\\
Macro-piece counts (\S\ref{sec:macro}) & 27 further tile counts (e.g.\ $3864$, $9940$, $62125$, $279174$) & \PL{}, \CE{} data\\
Conjecture C (\S\ref{sec:conjC}) & the set of tile counts is $\bigcup_k k^2\cdot\{\text{families}\}$ & existence half \PL{} on the proved ranges; exclusion half \CJ{}\\
Census (\S\ref{sec:census}) & $199\,140$ / $129\,108$ / $63\,826$ (non-squarefree $N\le10^6$) & \CE{} with label columns\\
$W^{3+}$ evidence (\S\ref{sec:w3}) & all one-leg words refuted for $(4,7)$, $(5,8)$ at $g=3$ & \CE{} (not formal)\\
\bottomrule
\end{tabular}
\end{center}

\subsection{Reading guide}
Section~\ref{sec:setting} defines the tile, the targets and the two tools used throughout (the ladder subdivision and the square closure).
Section~\ref{sec:thmW} states Theorem~W with a worked example and a proof sketch; the complete proof is Appendix~\ref{app:W}.
Sections~\ref{sec:collar} and~\ref{sec:macro} give the other two families of constructions.
Section~\ref{sec:conjC} states Conjecture~C and what is proved about it.
Section~\ref{sec:census} gives the census; Section~\ref{sec:w3} the computer evidence on one open family.
Section~\ref{sec:corr} lists corrections, Section~\ref{sec:notproved} lists what is not proved, and Section~\ref{sec:credits} gives credits with dates.
A reader who wants only the new mathematics can read Sections~\ref{sec:setting}, \ref{sec:thmW} and \ref{sec:conjC}.

\subsection{Credit and the use of an AI system}
The constructions in Sections~\ref{sec:thmW}--\ref{sec:macro} build on those of D.~Paliy (the $W/\beta$ framework, caps and collars, the $\alpha$-isosceles up-set, seeds), M.~Beeson and G.~Ro\c{s}u (the targets $W$, $T$ and the methods), B.~Herdt, M.~Laczkovich, Y.~X.~Zhang and J.~P.~Harries; Section~\ref{sec:credits} says exactly which statement is whose, with dates.
These results were obtained with the help of an AI system (Claude, Anthropic); the questions and the directions were the author's, and the proofs have been checked in Lean or by hand as labelled; the provenance table is in the repository README.

%=====================================================================
\section{Setting}\label{sec:setting}

\subsection{The tile}\label{sub:tile}
Throughout, $e<f$ are coprime positive integers and
\[
a=ef,\qquad b=f^2-e^2,\qquad c=f^2,\qquad \kappa=b+c=2f^2-e^2 .
\]
The \emph{tile} $\tile(e,f)$ is the triangle with sides $a,b,c$; $\alpha,\beta,\gamma$ are its angles opposite $a,b,c$.
Heron's formula gives $16\,\mathrm{Area}^2=e^2b^2(4f^2-e^2)$, and one checks $c^2=a^2+bc$ and
\[
3\alpha+2\beta=\pi,\qquad \gamma=2\alpha+\beta,\qquad \cos\alpha=\frac{\kappa}{2c}.
\]
These tiles are the tiles of type~(T) in the classification of the earlier paper: they are the tiles for which the prime case reduces to a question about $T_1(e,f)$.
The angle $\alpha/\pi$ is irrational; this is not used below.
Small examples: $(e,f)=(1,2)$ gives the tile $(2,3,4)$; $(2,3)$ gives $(6,5,9)$; $(2,5)$ gives $(10,21,25)$; $(1,4)$ gives $(4,15,16)$.
\begin{figure}[ht]
\FIGSLOT{tile}
\caption{The tile $\tile(e,f)$ with sides $a=ef$, $b=f^2-e^2$, $c=f^2$ and angles $\alpha$ (opposite $a$), $\beta$ (opposite $b$), $\gamma=2\alpha+\beta$ (opposite $c$), drawn for $(e,f)=(2,5)$, the tile $(10,21,25)$. The relation $3\alpha+2\beta=\pi$ is what makes the angles of the tile combine into the angles of the targets below.}
\end{figure}

\subsection{The targets}\label{sub:targets}
Let $M\ge1$ be an integer. The three families of targets that occur are listed in the table; the last column is the number of tiles that fit by area, which is the only count a tiling of the target can have.
\begin{center}
\small
\begin{tabular}{@{}llp{5.2cm}l@{}}
\toprule
target & sides & angles & area in tiles\\
\midrule
$T_M(e,f)$ & $Mf^3,\ Mf^3,\ Me(3f^2-e^2)$ & isosceles; apex $3\alpha$, base angles $\beta$ & $M^2(3f^2-e^2)$\\
$W_M(e,f)$ & $Mfb,\ Me\kappa,\ Mf^3$ & angles $2\alpha,\ \beta,\ \alpha+\beta$ (opposite $Me\kappa,\ Mfb,\ Mf^3$) & $M^2(2f^2-e^2)$\\
$B_g(e,f)$ & $gbc,\ gbc,\ gb\kappa$ & isosceles; base angles $\alpha$, apex $\alpha+2\beta$ & $b\kappa\,g^2$\\
\bottomrule
\end{tabular}
\end{center}
(For $W_M$, the side opposite the angle $2\alpha$ is $Me\kappa$, opposite $\beta$ is $Mfb$, and opposite $\alpha+\beta$ is $Mf^3$.)
$T_M$ and $W_M$ are the targets of Beeson--Ro\c{s}u and Paliy (``base-$\beta$'' targets); $B_g$ is the ``base-$\alpha$'' target (Paliy's $\alpha$-isosceles).
We call a target \emph{tiled} if it is tiled by copies of the tile $\tile(e,f)$ of the same $(e,f)$.
$W_M(e,f)$ and $T_M(e,f)$ are related by gluing: a tile at scale $Mf$ glued to the side $Mfb$ of $W_M$ gives $T_M$ (\S\ref{sub:Tcor}).
\begin{figure}[ht]
\FIGSLOT{targets}
\caption{The three targets for $(e,f)=(2,5)$ and $M=1$, drawn to scale: $W_1(2,5)$ (sides $105,92,125$), $T_1(2,5)$ (legs $125$, base $142$), and $B_1(2,5)$ (legs $525$, base $966$). The tile $(10,21,25)$ is drawn inside each for comparison. The angles are marked; they are built from $\alpha,\beta$ by $3\alpha+2\beta=\pi$.}
\end{figure}

\subsection{Two tools}\label{sub:tools}
\begin{proposition}[Ladder subdivision; \PL{}]\label{prop:ladder}
Let $s\ge1$ be an integer. A triangle similar to $\tile$ with sides $sa,sb,sc$ is tiled by $s^2$ copies of $\tile$.
A parallelogram with sides $sa$ and $uc$ enclosing the angle $\beta$ is tiled by $2su$ copies of $\tile$ (cut each unit cell along its diagonal of length $b$).
\end{proposition}
\emph{Lean:} \lean{Subdivision.ladderDissection}, \lean{Ladder.ladder}; the earlier paper calls this \texttt{thm:ladder}.
A \emph{trapezoid} below means a triangle of scale $s_1$ minus a corner triangle of scale $s_2<s_1$ of the same shape; it contains $s_1^2-s_2^2$ tiles.
\begin{proposition}[Square closure; \PL{}]\label{prop:square}
If $N$ is a tile count then so is $k^2N$ for every integer $k\ge1$.
\end{proposition}
\emph{Lean:} \lean{TileCountFamilies.tileCount_scale}. (Cut every tile of a tiling of $\Omega$ into $k^2$ triangles similar to it, with ratio $1/k$; the $k^2N$ pieces are congruent.)
\begin{remark}
In words: scaling the whole figure by $k$ multiplies the number of tiles by $k^2$.
Because of Proposition~\ref{prop:square} a tile count is naturally written $N=k^2n$, and thresholds below are thresholds in the \emph{scale of the target relative to the tile}, not in $k$.
\end{remark}


%=====================================================================
\section{Theorem W: the six-piece template}\label{sec:thmW}

\subsection{Statement}
Recall $a=ef$, $b=f^2-e^2$, $c=f^2$, $\kappa=2f^2-e^2$, and put $g=f-e$.
\begin{theorem}[Theorem W; \PL{} and \PH{}]\label{thm:W}
Let $e,f$ be coprime integers with $2\le e\le f-2$, and let $M\ge1$ be an integer with
\[
\text{(H1)}\quad Me\ge f\qquad\text{and}\qquad \text{(H2)}\quad (M-1)(f-e)\ge f .
\]
Then $W_M(e,f)$ is tiled by $M^2(2f^2-e^2)$ copies of $\tile(e,f)$ and $T_M(e,f)$ is tiled by $M^2(3f^2-e^2)$ copies of $\tile(e,f)$.
\end{theorem}
\emph{Lean:} \lean{Draft.ThmW.thmW}\,: for natural numbers \texttt{e f M} with \texttt{Nat.Coprime e f}, \texttt{2 <= e}, \texttt{e+2 <= f}, \texttt{f <= M*e}, \texttt{f <= (M-1)*(f-e)}, the statement \texttt{WTilable e f M $\wedge$ TileCount ((2*f\^{}2-e\^{}2)*M\^{}2) $\wedge$ TileCount (M\^{}2*(3*f\^{}2-e\^{}2))}.
Here \texttt{WTilable e f M} says that there is a dissection into congruent triangles, each congruent to the triangle with sides $ef,\ f^2-e^2,\ f^2$, of a triangle with sides $Mf(f^2-e^2)$, $Me(2f^2-e^2)$, $Mf^3$; the two \texttt{TileCount} conjuncts say only that the number is a tile count (the shape is pinned by the first conjunct).
The module is \texttt{Draft/ThmW.lean} with its \texttt{ThmW*} helper modules; the aggregate is \lean{Erdos634.AllThmW}; the axioms are \texttt{propext}, \texttt{Classical.choice}, \texttt{Quot.sound}.
The statement of the theorem was compared with the hand statement in an audit inside the project (verdict: match, with a note about credits only).
The hand proof, given in full in Appendix~\ref{app:W}, was audited separately by an adversarial reader and by exact computation (Section~\ref{sub:Wevidence}).

\why{Before this theorem, $W_M(e,f)$ and $T_M(e,f)$ were known to be tiled for $M\ge f$ (Paliy, Beeson), and the working expectation was that no smaller $M$ works for $e\ge2$; the only smaller scales known were the cases $e=1$ and $e=f-1$, where the threshold is indeed $f$.
Theorem~W shows that for $e\ge2$ the scale $f$ is \emph{not} the threshold: for $(e,f)=(2,5)$ the scale $3$ works instead of $5$, and for $e<f/2$ the new threshold is exactly the proved lower bound (Corollary~\ref{cor:sharp}).}

\medskip
Put $M^*(e,f)=\max\bigl(\lceil f/e\rceil,\ \lceil f/(f-e)\rceil+1\bigr)$; the hypotheses (H1), (H2) say $M\ge M^*(e,f)$.
\begin{center}\small
\begin{tabular}{@{}lllll@{}}
\toprule
$(e,f)$ & $M^*$ & first tiled target & number of tiles & previously known threshold\\
\midrule
$(2,5)$ & $3$ & $W_3(2,5)$ & $414$ & $5$\\
$(3,5)$ & $4$ & $W_4(3,5)$ & $656$ & $5$\\
$(3,7)$ & $3$ & $W_3(3,7)$ & $801$ & $7$\\
$(2,11)$ & $6$ & $W_6(2,11)$ & $8568$ & $11$\\
$(9,11)$ & $7$ & $W_7(9,11)$ & $7889$ & $11$\\
\bottomrule
\end{tabular}
\end{center}
The thirteen values $M<f$ with $4f<9e$ that were found first by a computer search are listed in Appendix~\ref{app:data}; they are instances of the theorem.

\subsection{Consequences}
\begin{corollary}[Sharpness for $e<f/2$; \PL{}]\label{cor:sharp}
Let $2\le e<f/2$, $\gcd(e,f)=1$. Then $W_M(e,f)$ is tiled by copies of $\tile(e,f)$ if and only if $Me\ge f$.
\end{corollary}
\begin{proof}
If $Me<f$ there is no tiling: this is the earlier \PL{} corollary \texttt{cor:ebelow}(a) (Lean \lean{Draft.SideRunG.ebelow_a}, in the release paper \cite{ReleasePaper}).
Conversely, $e<f/2$ gives $f-e>f/2$, so $f/(f-e)<2$ and $\lceil f/(f-e)\rceil+1=3$, while $f/e>2$ gives $\lceil f/e\rceil\ge3$; hence $M^*(e,f)=\lceil f/e\rceil$ and (H1) implies (H2). Apply Theorem~\ref{thm:W}.
\end{proof}
\begin{corollary}[Answer to an open case; \PL{}]\label{cor:W435}
$W_4(3,5)$, with sides $320,\ 492,\ 500$, is tiled by $656$ copies of $(15,16,25)$. A tiling is shipped as exact data (file \texttt{tiling\_656\_W\_3\_5\_M4}).
\end{corollary}
\noindent The earlier paper left the scales $M<f$ open apart from a few instances, and recorded the pattern ``every case decided so far is tiled exactly when $M\ge f$'' (\cite{ReleasePaper}, Section on scales).
Corollary~\ref{cor:W435} is a counterexample to that pattern: $M=4<5=f$.
It also refutes a size conjecture of the working notes (a target not similar to the tile has longest side at least $c^2$): here $500<625=c^2$ (see Section~\ref{sec:corr}).

\subsection{The $T$ half}\label{sub:Tcor}
Glue to the side $Mfb$ of $W_M(e,f)$ one copy of the tile at scale $Mf$ (sides $Mef^2$, $Mfb$, $Mf^3$) with its side $Mfb$ on that side and its far vertex $Z$ on the line $AB$, beyond $B$ (so $BZ=Mef^2$); the result is the isosceles triangle with legs $Mf^3$ and base $Me(3f^2-e^2)$, that is $T_M(e,f)$, and it has $M^2(2f^2-e^2)+M^2f^2=M^2(3f^2-e^2)$ tiles.
This is the earlier \texttt{prop:layerall}(a) (Lean \lean{CollarG.betaTilable_of_wTilable}); the Lean statement of Theorem~W includes it (\lean{DebtB.counts_of_wTilable} needs only the weaker hypothesis $M(f-e)\ge f$ on the $T$ side).
The converse (a tiling of $T_M$ gives one of $W_M$) is not known; the proved exclusions for $T_M$ are weaker than for $W_M$ (only $2Me<f$, and $e=1$, $2M\le f$), so the $T$ column of the census has more open scales (Section~\ref{sec:notproved}).

\subsection{The six pieces, in words, and a worked example}\label{sub:six}
\textbf{Idea.} The target $W_M(e,f)$ has a vertex $O$ with angle $2\alpha$ and two other vertices $A$ (angle $\beta$) and $B$ (angle $\alpha+\beta$).
Let $P$ be the point of the side $AB$ with $OP=OB$; then $OP$ bisects the angle at $O$, and the triangle $OBP$ is isosceles with apex angle $\alpha$.
The template cuts $W_M$ into six pieces, each of which is tiled by Proposition~\ref{prop:ladder}:
\begin{center}\small
\begin{tabular}{@{}llp{6.6cm}l@{}}
\toprule
piece & kind & scales (integers) & tiles\\
\midrule
$\pi_0$ & trapezoid & triangle $OAP$ of scale $fM$ minus a corner triangle of scale $et$ & $f^2M^2-e^2t^2$\\
$\pi_1$ & triangle & scale $v=fM-er$ & $v^2$\\
$\pi_2$ & triangle & scale $ft$ & $f^2t^2$\\
$\pi_3$ & trapezoid & scales $s_1=f(r-t)$ and $s_2=fr-eM$ & $s_1^2-s_2^2$\\
$\pi_4$ & parallelogram & sides $a s_1$ and $c\,x$, $x=ft-eq$, angle $\beta$ & $2s_1x$\\
$\pi_5$ & trapezoid & scales $u_1=fq-et$ and $u_2=fq-er$ & $u_1^2-u_2^2$\\
\bottomrule
\end{tabular}
\end{center}
with four free integers $M,t,r,q$.
The six tile counts add up to $M^2(2f^2-e^2)$ identically in $e,f,M,t,r,q$ (a polynomial identity), so what remains is to place the pieces: all vertices are given by vector sums with coefficients linear in $M,t,r,q$, and ``the pieces lie in the target and do not overlap'' becomes a list of $26$ inequalities between linear and quadratic expressions in $e,f,M,t,r,q$ (Appendix~\ref{app:W}, Lemma~B).
The proof chooses
\[
n=\Bigl\lfloor \frac{eM}{f}\Bigr\rfloor,\quad \rho=eM-fn,\quad \delta=\begin{cases}1&\rho=0\\0&\rho>0\end{cases},\qquad t=n-\delta,\quad r=n+1,\quad q=n ,
\]
and shows that under (H1), (H2) every inequality holds (Lemma~C).
\begin{example}[$W_3(2,5)$ with $414$ tiles]\label{ex:414}
$(e,f,M)=(2,5,3)$: $a=10$, $b=21$, $c=25$, $\kappa=46$, and $W_3(2,5)$ has sides $315,\ 276,\ 375$.
Here $n=\lfloor6/5\rfloor=1$, $\rho=1$, $\delta=0$, so $(t,r,q)=(1,2,1)$, and
\begin{center}\small
\begin{tabular}{@{}lccccc@{}}
\toprule
 & scales & tiles & & scales & tiles\\
\midrule
$\pi_0$ & $15$ minus $2$ & $225-4=221$ & $\pi_3$ & $5$ minus $4$ & $25-16=9$\\
$\pi_1$ & $v=11$ & $121$ & $\pi_4$ & $a s_1\times c x$, $s_1=5,\ x=3$ & $2\cdot5\cdot3=30$\\
$\pi_2$ & $5$ & $25$ & $\pi_5$ & $3$ minus $1$ & $9-1=8$\\
\bottomrule
\end{tabular}
\end{center}
The total is $221+121+25+9+30+8=414=9\cdot46$. A full list of the $414$ tiles is shipped (file \texttt{tiling\_414\_W\_2\_5\_M3}) and has been verified by an independent exact checker.
\end{example}
\begin{figure}[ht]
\FIGSLOT{sixpiece}
\caption{The six-piece layout of Theorem~W for $W_3(2,5)$, a triangle with sides $315,\ 276,\ 375$ tiled by $414$ copies of $(10,21,25)$. The pieces $\pi_0,\dots,\pi_5$ are shaded differently and labelled with their kind and scales as in Example~\ref{ex:414}: $\pi_0$ is the large trapezoid at the vertex $O$ of angle $2\alpha$ (triangle $OAP$ of scale $15$ minus a corner of scale $2$), $\pi_1$ the triangle of scale $11$ on the side $OB$, $\pi_2$ the triangle of scale $5$ on the side $AB$, $\pi_3$ and $\pi_5$ the two small trapezoids, $\pi_4$ the parallelogram of $30$ tiles. The point $P$ on $AB$ with $OP=OB$ is marked.}
\end{figure}

\subsection{Proof sketch}\label{sub:Wsketch}
(1)~\emph{Shapes.} Each piece is the stated union of tiles when its scales are positive integers: the triangles by SAS/SSS, the trapezoids because the corner triangle is cut off by a line parallel to a side, the parallelogram by Proposition~\ref{prop:ladder}. The shape claims are identities of Eisenstein-type numbers $x+i\,w\,y$ with $w=\sqrt{4f^2-e^2}$ and use only $3\alpha+2\beta=\pi$ (equivalently $\varepsilon^3\eta^2=-1$ for $\varepsilon=e^{i\alpha}$, $\eta=e^{i\beta}$) and the closure relation of the tile.
(2)~\emph{Placement.} The pieces are convex and counter-clockwise; they lie in the target iff their vertices do, and two pieces have disjoint interiors if one edge line of one piece separates it from the other. Together with the area identity this gives a tiling. Each needed sign is a product of positive constants and of factors from a fixed list of $26$ linear or quadratic expressions in $e,f,M,t,r,q$.
(3)~\emph{Choice of parameters.} With $n,\rho,\delta,t,r,q$ as above and $s=M-n$ (so $gM=fs-\rho$ for $g=f-e$), every factor is positive. The only places where (H1) and (H2) are used are: $n\ge1$ and $gn\ge e+1$ (from (H1) and (H2) and coprimality), and $s\ge2$ (from (H2)); the decisive inequality is the identity $\text{(LB)}=s\,\Psi-\kappa(\rho+e)$ with $\Psi=f^3+2ef^2-e^2f-e^3>0$, which is positive because $s\ge2$ and $\rho\le f-1$.
(4)~\emph{Lean.} Items (1)--(3) are formalised as separate modules (point lemmas, shape lemmas, one sign lemma per factor, the parameter lemma, the packing); the packing uses the general theorem that a finite family of convex pieces, each tiled, with disjoint interiors and total area equal to the area of the target, tiles it.

\subsection{Evidence besides the proof}\label{sub:Wevidence}
\CE{} For \emph{every} coprime pair $2\le e\le f-2$ with $f\le300$ (26\,800 pairs), an exact integer program checked the scale positivity, the SSS shapes, the trapezoid ratios, convexity, containment, the area sum and a separating edge for all $15$ pairs of pieces, for $M$ from $M^*$ to $M^*+7$ and for $M\in\{f,f+1,2f-1,2f,3f\}$: $348\,350$ cases, no failure. Negative controls: for $M=M^*-1$ the canonical parameters fail for every pair tried. A second, independent checker (own code, a different arithmetic) repeated this for all $f\le120$ ($706\,489$ cases) and for random $f\le400$.
Thirty-eight explicit tilings were expanded to the full list of tiles and verified by an exact checker (congruence, pairwise disjointness, area, the hull is a triangle), among them $W_3(2,5)$, $W_3(3,7)$, $W_5(2,9)$ ($3950$ tiles), $W_8(11,13)$ ($13\,888$ tiles), $W_4(11,43)$ ($57\,232$ tiles) and $T_3(2,5)$, $T_6(3,16)$, $T_7(24,29)$ ($95\,403$ tiles).
No instance of the theorem contradicts a proved exclusion (checked for all $e,f\le200$: $4\,786\,740$ admissible triples), and none of the $197$ recorded exclusions of the census registry meets the family.

\subsection{What stays open}
For $e<f/2$ nothing: Corollary~\ref{cor:sharp}.
For $e>f/2$ the \emph{no} side of (H2) is open: the proved exclusion is only $M(f-e)\le e$ (earlier \texttt{cor:thickw}(a)), so for each row $(e,f)$ exactly the (at most two) scales $M$ with $(M-1)(f-e)<f<(M+1)(f-e)$ are undecided; the conjecture is that $W_M$ is tiled iff $(M-1)(f-e)\ge f$.
The members $e=1$ and $e=f-1$ are outside the theorem; for both the threshold is known to be $M=f$ (\texttt{cor:below}, \texttt{cor:thickw}(b)).
For $T_M$ the open scales are those with $2Me\ge f$ and $(M-1)(f-e)<f$, the first being $T_3(1,4)$ with $e=1$ and, for $e\ge2$, for example $T_2(3,5)$ and $T_3(3,5)$.


%=====================================================================
\section{The collar family and the base-$\alpha$ rows}\label{sec:collar}

This section is about the third target of Section~\ref{sub:targets}, the isosceles triangle $B_g(e,f)$ with base angles $\alpha$, legs $gbc$, base $gb\kappa$, and $b\kappa g^2$ tiles.
Write $N_0(e,f)=b\kappa=(f^2-e^2)(2f^2-e^2)$, so that a tiling of $B_g(e,f)$ has $N_0(e,f)g^2$ tiles; \emph{the row $(e,f)$} is the set of $g$ for which $B_g(e,f)$ is tiled.
The first values: $N_0(1,2)=21$, $N_0(2,3)=70$, $N_0(1,4)=465$, $N_0(2,5)=966$, $N_0(1,6)=2485$.
The row is non-empty in every case examined, and it is closed under addition (Paliy, \CP{}; his construction), so it is, if his claim holds, a numerical semigroup; what is open is its smallest element and its gaps.

\subsection{A priority statement}
D.~Paliy stated on 29 September 2026 \cite{PaliyRepo} that the tile counts of the form $b\kappa K^2$ are exactly the counts of tilings of the targets $B_K$ and that, for each tile, every $K$ beyond an explicit threshold $K_0$ is realised (the ``$\alpha$ up-set''); the seeds $84$ and $189$ of the row $(1,2)$ are in Beeson's work (the tilings of $84$ and $189$ in his Figures 28 and 29; $189$ is due to Harries).
\textbf{The up-set and the counts are Paliy's.}
What is ours in this section is an explicit construction of the layers (the \emph{collar}), a parametric seed for the whole family $e=1$, the explicit thresholds that follow, and their Lean proofs; the construction is independent of his, and the thresholds below are lower than the ones obtained from his formula as we read it for the rows $(1,4)$, $(1,6)$ and $(2,5)$ of Theorem~\ref{thm:tails} (we have not asked him).

\subsection{The collar step}
\begin{theorem}[Collar step; \PH{}, with a \PL{} one-layer case]\label{thm:collar}
Let $e<f$ be coprime, $a=ef$, $b=f^2-e^2$, $\kappa=2f^2-e^2$, and suppose $B_{g_0}(e,f)$ is tiled. Choose $s\in[0,a)$ with $bs\equiv b\kappa g_0\pmod a$ (it exists since $\gcd(a,b)=1$), put $K=(b\kappa g_0-bs)/a$ and $h_{\max}=\lfloor K/a\rfloor$.
Then for every integer $h$ with $1\le h\le h_{\max}$ the triangle $B_{g_0+h}(e,f)$ is tiled.
\end{theorem}
\emph{Construction.} $B_{g_0+h}$ is a homothetic copy of $B_{g_0}$ at the apex plus a trapezoidal strip along the base; the strip is cut into three pieces by two lines, each tiled by Proposition~\ref{prop:ladder}: a triangle of scale $ch$, the difference of two triangles $\Delta(K)\setminus\Delta(K-ah)$, and the difference $\Delta(s+bh)\setminus\Delta(s)$.
The identity that makes the cut edges coincide is $c^2h+a(K-ah)=b\kappa(g_0+h)-b(s+bh)$, which uses only $c^2=a^2+bc$ and $\gamma=2\alpha+\beta$; the range $h\le K/a$ is exactly the condition $K-ah\ge0$.
\emph{Status.} The hand proof was audited independently (verdict: sound; checked exactly on $17\,853$ instances, all coprime $e<f<40$ with several $(g_0,h)$ each, and no failure). \PL{} The case $h=1$, for every row with $ef\le f^2-e^2$, is Lean \lean{Draft.CollarFam.layer_step} (real-level core \lean{Draft.CollarLayer.layer_core}); arbitrary $h$ is hand.
The corollary ``$g_0$ in the row implies every $g\ge g_0$ is in the row'' needs $h_{\max}(g_0)\ge g_0-1$ together with $h_{\max}(x)\ge1$ for $x\ge2g_0-1$ (the first version of this claim in our notes omitted the second condition and was corrected by the audit); it holds for the rows below.

\begin{figure}[ht]
\FIGSLOT{collar}
\caption{The collar: $B_{g_0+h}(1,4)$ with $g_0=5$, $h=3$ (the tiling of $465\cdot8^2=29\,760$ tiles). The inner triangle is the homothetic copy of $B_{5}(1,4)$ (tiled by $11\,625$ tiles, itself a six-piece layout); the strip below it is cut into the three pieces of Theorem~\ref{thm:collar}, labelled with their scales. Each piece is a ladder subdivision.}
\end{figure}

\subsection{The seed for $e=1$}
\begin{theorem}[Seed and up-set for $e=1$; \PL{}]\label{thm:seedA}
For every $f\ge2$ and every $g\ge f+1$, $\,(f^2-1)(2f^2-1)\,g^2$ is a tile count; the tiling is of $B_g(1,f)$ by copies of $(f,f^2-1,f^2)$.
\end{theorem}
\emph{Lean:} \lean{Draft.CollarUp.tileCount_e1_up}. The seed $g=f+1$ is a six-piece layout whose pieces are polynomial in $f$ (generated and checked symbolically; module \texttt{CollarSeed}, $125$ lemmas), and the layer step then gives all $g\ge f+1$.
For $f=2,3,4,6$ the seed counts are $189$, $2176$, $11\,625$, $121\,765$; the first is the $189$ of Beeson and Harries, the third and fourth were found by a template search in release v1.16.
\why{The sporadic-looking counts $189$, $11\,625$, $121\,765$ of earlier releases are the first members of one family, and the count $(f^2-1)(2f^2-1)g^2$ is a tile count for all $f\ge2$ at scales that no longer depend on a search.}

\subsection{Lowered tails for three rows}
\begin{theorem}[Tails; \PL{}]\label{thm:tails}
(a) $465\,g^2$ is a tile count for every $g\ge5$ (row $(1,4)$, tile $(4,15,16)$).
(b) $2485\,g^2$ is a tile count for every $g\ge2$ (row $(1,6)$, tile $(6,35,36)$).
(c) $966\,g^2$ is a tile count for every $g\ge2$ (row $(2,5)$, tile $(10,21,25)$).
\end{theorem}
\emph{Lean:} (a) \lean{Draft.CollarFam.tileCount_465_sq}, from the seed $11\,625$ and the layer step. (b) The tail $g\ge5$ is \lean{Draft.CollarFam5.tileCount_2485_sq5}; $g=2,3$ are $9940$ and $22\,365$ (Section~\ref{sec:macro}), and $g=4$ is $4\cdot9940$ by Proposition~\ref{prop:square}. (c) The tail $g\ge17$ is \lean{Draft.CollarFam17.tileCount_966_sq17}; $g=2,\dots,8$ and $10,\dots,18$ are the macro-piece counts $3864,\dots,312\,984$ of Section~\ref{sec:macro} (sixteen values), and $g=9$ is $9\cdot8694=78\,246$ (\lean{Draft.AllMacro.tileCount_78246}).
The squarefree cores $465$, $966$ and $2485$ themselves ($g=1$) are not tile counts known to us; they are among the open squarefree values.
\begin{remark}
In our working notes we had recorded the pattern ``the smallest scale of a row is $f+1$ for even $f$''. (c) and (b) refute it: the rows $(2,5)$ and $(1,6)$ contain $g=2$.
For $(1,4)$ the values $g=2,3,4$ ($1860$, $4185$, $7440$) are not known to be tile counts or non-tile counts.
Paliy's threshold, read from his formula, is $K_0=17$ for $(2,5)$ (equal to our tail in (c) before the macro pieces fill $g=2,\dots,16$), $6$ for $(1,4)$ and $8$ for $(1,6)$ (we did not check these values with him).
\end{remark}

%=====================================================================
\section{Macro-piece tile counts}\label{sec:macro}

\subsection{The idea}
In the six-piece layout of Section~\ref{sec:thmW} every piece is tiled by the ladder subdivision, that is, by a tiling of a \emph{triangle similar to the tile}.
A \emph{macro piece} is a piece that is a triangle similar, with integer ratio $k$, to the \emph{target} of a tiling already known: such a piece is tiled by $k^2n$ tiles if the known tiling has $n$ tiles (Lean: \lean{Draft.MacroPiece.macroTFam_gen}, which transports a known dissection by the ladder theorem and congruence of triangles with equal sides).
Macro pieces were used for a long time by hand (Laczkovich's macro-dissections, Herdt's $1125$; Section~6 of the version-2.2 paper), and they enlarge the set of layouts for which the area equation has integer solutions.
The layouts are found by a computer search over the combinatorial type of the layout and the integer scales of its pieces, with the equations solved exactly; a layout is then expanded to the full list of tiles, verified by an exact checker, and finally turned into a Lean module, whose proof is the same area-and-containment argument as in Theorem~W with the numbers of that layout.
The search itself is not part of this release; what is shipped are the results: the exact tilings (ninety-seven files, Appendix~\ref{app:data}) and the Lean theorems.

\subsection{The counts}
\begin{theorem}[Macro-piece tile counts; \PL{}, \CE{} data]\label{thm:macro}
Each of the following is a tile count, as the statement \texttt{TileCount $N$} of a Lean theorem:
\begin{center}\small
\begin{tabular}{@{}p{4.3cm}p{2.5cm}p{2.5cm}p{3.1cm}@{}}
\toprule
values $N$ & form & tile & what is tiled\\
\midrule
$375$, $540$, $735$ & $15m^2$, $m=5,6,7$ & $(3,5,7)$ & equilateral triangles of side $75,90,105$\\
$825$ & $33\cdot5^2$ & $(3,5,7)$ & isosceles, sides $105,105,165$\\
$630$, $1120$, $1750$ & $70\,g^2$, $g=3,4,5$ & $(6,5,9)$, row $(2,3)$ & $B_g(2,3)$\\
$3864$, $8694$, $15\,456$, $24\,150$, $34\,776$, $47\,334$, $61\,824$ & $966\,g^2$, $g=2,\dots,8$ & $(10,21,25)$, row $(2,5)$ & $B_g(2,5)$\\
$96\,600$, $116\,886$, $139\,104$, $163\,254$, $189\,336$, $217\,350$, $247\,296$, $279\,174$, $312\,984$ & $966\,g^2$, $g=10,\dots,18$ & $(10,21,25)$ & $B_g(2,5)$\\
$9940$, $22\,365$, $62\,125$ & $2485\,g^2$, $g=2,3,5$ & $(6,35,36)$, row $(1,6)$ & $B_g(1,6)$\\
\bottomrule
\end{tabular}
\end{center}
That is $26$ values, and a twenty-seventh, $78\,246=9\cdot8694$, by Proposition~\ref{prop:square}.
\end{theorem}
\emph{Lean:} one module \texttt{Draft/TM$N$} (or \texttt{Draft/T62125}) per value, with \texttt{theorem tileCount\_$N$ : TileCount $N$}; aggregate \lean{Draft.AllMacro}. Each module fixes the layout, proves the shapes and the containment by the area-and-sign method, and does not use a tile list; the standard axioms only.
The exact tilings of all of them are in \lean{data/witnesses/v3_exact} (for $N=375$, \texttt{tiling\_375\_macro\_5\_3\_7\_g1}); they have been verified by an exact checker independently of the Lean proofs.
\why{Together with Theorem~\ref{thm:tails} these counts show that in two rows, $(2,5)$ and $(1,6)$, \emph{every} scale $g\ge2$ is realised; the squarefree cores of those rows are the only open members. The $15m^2$ values for $m=5,6,7$ are the seeds that Paliy used (\CP{}) for his claim that $15m^2$ is a tile count for every $m\ge4$.}
For the equilateral counts $15m^2$: $m=4$ ($240$) was already a Lean theorem, and so was every $m$ in the semigroup generated by $4$ and $9$ (earlier release).
The Lean theorem \lean{Draft.EqSemigroup.eqD_step} says that equilateral tilings of sides $15m_1$ and $15m_2$ by copies of $(3,5,7)$ give one of side $15(m_1+m_2)$; the seeds $4,5,6,7$ would therefore give every $m\ge4$ (this is Paliy's claim, \CP{}).
The new Lean theorems for $m=5,6,7$ state only \texttt{TileCount} and are not yet instances of the equilateral predicate used by \texttt{eqD\_step}, so ``$15m^2$ for all $m\ge4$'' is \emph{not} a Lean theorem here; the exact tilings (sides $75,90,105$) are verified, and $15m^2$ for the individual $m=5,6,7$ is Lean.

\begin{figure}[ht]
\FIGSLOT{macro}
\caption{Composition of a macro-piece tiling: a template layout of $B_g(2,5)$ with $g=6$ ($34\,776=966\cdot36$ tiles) whose pieces are ladder triangles, parallelograms, trapezoids and one \emph{macro piece}, a triangle similar with integer ratio to the target of the already-known tiling of $966\cdot g'^2$ tiles; the macro piece is drawn with the known tiling inside it.}
\end{figure}

\subsection{Other counts shipped in this release}
Twelve counts that are squares times earlier template counts ($46\,500$, $104\,625$, $267\,840$, $290\,625$, $418\,500$, $476\,160$, $487\,060$, $569\,625$, $602\,640$, $671\,460$, $744\,000$, $941\,625$; module \texttt{Draft/TemplateSquares}) are Lean theorems by Proposition~\ref{prop:square}; they were already tile counts mathematically and had been omitted by the census bookkeeping.
The thirteen tilings of $W_M$ with $M<f$ of Appendix~\ref{app:data} are also Lean theorems (\texttt{Draft/AllWBelow}), as are their upward closures (\texttt{Draft/AllDebtB}); they are subsumed by Theorem~W but were proved first.
$4830$ is a tile count: Paliy's certificate (4830 copies of $(24,11,31)$ tile the triangle $(961,1155,1426)$) has been converted and verified exactly (\CE{}; the construction is Paliy's, \CP{}; his family of mixed-corner layouts gives further values that we have not verified).


%=====================================================================
\section{Conjecture C: a candidate classification of the tile counts}\label{sec:conjC}

The earlier papers give the exclusion side and a growing list of tile counts, but no statement of what the set of tile counts \emph{is}.
This section states one, as a conjecture, and says exactly what part of it is proved.
\textbf{It is not claimed to be true.} It is an organising map: a new tile count is either in a family (no news), or extends a family's range, or is a counterexample to the conjecture.

\subsection{The statement}
Recall Proposition~\ref{prop:square}: every family below is understood closed under $n\mapsto k^2n$.
\begin{conjecture}[Conjecture C; \CJ{}]\label{conj:C}
A positive integer $N$ is a tile count if and only if $N=k^2n$ for some integer $k\ge1$ and some core $n$ in one of the families $F_0,\dots,F_8$ of Table~\ref{tab:families}, with the parameter ranges stated there.
\end{conjecture}
\begin{table}[ht]
\small
\renewcommand{\arraystretch}{1.15}
\begin{tabular}{@{}p{0.55cm}p{3.2cm}p{3.3cm}p{5.1cm}@{}}
\toprule
 & cores $n$ & tile, target & range proved \quad (label) / credit\\
\midrule
$F_0$ & $1$, $x^2+y^2$ ($x,y\ge1$), $3$, $6$ & similar tiles, rep-tiles & all \quad \PL{} / Snover--Waiveris--Williams, Laczkovich\\
$F_1$ & $2f^2-e^2$ and $3f^2-e^2$, at scale $M$: $M^2(2f^2-e^2)$, $M^2(3f^2-e^2)$ & $\tile(e,f)$; $W_M$, $T_M$ & $e=1$ or $e=f-1$: $M\ge f$. $2\le e\le f-2$: $M\ge M^*(e,f)$ (Thm~\ref{thm:W}). \quad \PL{} / Paliy, Beeson ($M\ge f$); ours ($M^*$)\\
$F_2$ & $(2q^2-p^2)(3q^2-p^2)$, $0<p<q$ & five-piece layouts & all \quad \PL{} / Paliy\\
$F_3$ & $(2a+b)(a+b)$, $3(a+b)(a+2b)$, $(a+2b)(2a+b)$ and Zhang's five forms $\times m^2$; $15m^2$, $33m^2$ & tiles with an angle $2\pi/3$, $c^2=a^2+ab+b^2$; $(3,5,7)$ & ranges as in the Lean predicates; $15m^2$: $m\in\langle4,9\rangle$ and $m=5,6,7$ \quad \PL{} / Paliy, Zhang, Herdt. $15m^2$ for all $m\ge4$: \CP{}\\
$F_4$ & $88$, $143$ (sporadic) & tile $(3,5,7)$ on the triangles $(21,55,56)$ ($88$) and $(39,49,55)$ ($143$) & \PL{} / Harries ($88$); ours ($143$)\\
$F_5$ & $(f^2-e^2)(2f^2-e^2)\,g^2$, $g\in G(e,f)$ & $\tile(e,f)$; $B_g(e,f)$ & $e=1$: $g\ge f+1$; $(1,4)$: $g\ge5$; $(1,6),(2,5)$: $g\ge2$; $(2,3)$: $g=3,4,5$ and more by Thm~\ref{thm:collar} \quad \PL{} / \PH{}; Paliy: eventual up-set in every row, \CP{}\\
$F_6$ & $b\,(v(D+\tau v))^2$ with two Diophantine conditions; $1125$, $19\,600$, $130\,977$ & $\gamma=2\alpha$, tile $(u^2,v^2-u^2,uv)$ & Herdt's family and its uniform version \quad \PL{} / Herdt, Laczkovich\\
$F_7$ & $4830$ & $(24,11,31)$ & \CE{} / Paliy, \CP{} (mixed-corner family)\\
$F_8$ & $b\,t^2$, $t\in S(e,f)$ & $\tile(e,f)$; $\theta$-isosceles target (legs $tfb$, base $teb$, base angles $\alpha+\beta$, $bt^2$ tiles) & \CP{} / Paliy (semigroup $S(e,f)$, $te\ge2f$ necessary \PL{})\\
\bottomrule
\end{tabular}
\caption{The families of Conjecture~C and what is proved about each. $\langle4,9\rangle$ is the set of non-negative combinations of $4$ and $9$. In $F_1$ the ranges are those of Theorem~\ref{thm:W} and of the earlier theorem for $M\ge f$; for the members $e=1$ and $e=f-1$ the threshold $f$ is proved sharp for $W_M$.}\label{tab:families}
\end{table}
\emph{Remarks.} (i) $84=21\cdot2^2$ and $189=21\cdot3^2$ are not sporadic: they are the members $g=2,3$ of the row $(1,2)$ of $F_5$ (earlier lists had them in $F_4$); $77\in F_2$; $183$ is \emph{not} a tile count (it is among the excluded values $\le200$; an earlier draft of this section listed it by a misreading, see Section~\ref{sec:corr}).
(ii) Theorem~\ref{thm:W} is stated for $2\le e\le f-2$; $F_1$ in Table~\ref{tab:families} includes the two edge members with their threshold $f$.
(iii) For $F_6$ the earlier papers state two conjectures about its exclusion side, $V$ ($v\mid g$ is necessary) and $A'$ (a size bound); Section~\ref{sec:w3} gives evidence for them.

\begin{figure}[ht]
\FIGSLOT{t143}
\caption{The tiling of the triangle with sides $(39,49,55)$ by $143$ copies of the triangle $(3,5,7)$ (angle $2\pi/3$ between the sides $3$ and $5$), the sporadic count $143$ of $F_4$. It was found by a search in an earlier release and is a Lean theorem there; $143$ is not of any of the forms $k^2$, $x^2+y^2$, $3k^2$, $6k^2$ and is not in any parametric family of this paper.}
\end{figure}

\subsection{What is proved about C}
\begin{theorem}[The existence half, on the proved ranges; \PL{}]\label{thm:Cexist}
Let $\mathrm{ConjCBase}(N)$ be the statement ``$N=k^2n$, $k\ge1$, $n$ in the Lean predicate $F_0\cup\dots\cup F_6$ with the ranges proved at the time of that module''. Then $\mathrm{ConjCBase}(N)$ implies that $N$ is a tile count.
\end{theorem}
\emph{Lean:} \lean{ConjC.Exist.conjC_existence} (modules \texttt{ConjC/Defs}, \texttt{ConjC/Exist}, \texttt{ConjC/BetaCount}; aggregate \lean{AllConjC}; each family predicate is exactly the parameter range of an existing Lean existence theorem, and the non-vacuity of each is checked on an instance, e.g.\ \lean{conjCBase_143}, \lean{conjCBase_11625}). It collects constructions under one statement; it is not new mathematics beyond the family theorems it cites.
The predicate in Lean is narrower than Table~\ref{tab:families}: it has $F_1$ for $M\ge f$ only (Theorem~\ref{thm:W}, the tails and the macro counts of Sections~\ref{sec:collar}--\ref{sec:macro} are separate Lean theorems, not yet folded into it), it has no $F_7$, $F_8$, and its $F_5$ is the part proved at the time (the row $e=1$, $g\ge f+1$; $966g^2$, $g\ge25$; five engine values).

\medskip
\noindent\CE{} \emph{Consistency checks} (exact computations in the working notes, repeated for the census of this release where stated).
(a) The families generate exactly the $105$ tile counts $\le200$ of the earlier theorem, no more and no fewer.
(b) None of the $199\,140$ excluded values of the census is generated by any family (rechecked for the Theorem~W family and the macro counts in this release: no overlap).
(c) Every one of the known tile counts of the census is generated by some family; this check is partly circular, because the known values were found by these families, so it shows that the families are complete for what is known and nothing more.
The sharp content of C is the exclusion half.

\subsection{What the exclusion half needs}
Everything proved in the exclusion direction in the earlier releases is compatible with C. The exclusion half of C is not proved: it is the statement that no value outside the families is a tile count, and it decomposes into threshold statements, one for each family at the place where the known range stops:
\begin{itemize}
\item[$W^*$] for $F_1$: $W_M(e,f)$ is tiled iff $M\ge M^*(e,f)$; open for $e>f/2$ (Section~\ref{sec:thmW}); and $T^*$: $T_M$ is tiled iff $W_M$ is.
\item[$V,A'$] for $F_6$: the target $\Omega_g$ of $\gamma=2\alpha$ is tiled only if $v\mid g$, and for $g=1$, $(g-1)(v-u)<u$; open.
\item[$\Theta$] for $F_8$: $\theta_t(e,f)$ is tiled iff $t$ is in a semigroup, with $te<2f$ excluded (proved, \PL{}).
\item[$S_3$] for $F_5$: the row $G(e,f)$ has a smallest element $g_{\min}(e,f)$ and its gaps; known values are $g_{\min}(1,2)=2$ (exactly, since $21$ is not a tile count), $g_{\min}(2,5),\,g_{\min}(1,6)\in\{1,2\}$, $g_{\min}(1,4)\le5$.
\item[$Z$] for $F_3$: Zhang's counts $m^2a b$ etc.\ below the threshold $m<3t_0$ are not tile counts unless another family covers them; the known exceptions are the $15m^2$ and $33m^2$ counts.
\item small leaves: the $\pi/3$ and $2\pi/3$ pairs with irrational angle (closed, \PL{}), the commensurable cases.
\end{itemize}
Each of these is an infinite non-existence theorem; the methods of this project (additive invariants and the run induction of the earlier papers) prove them only in the ranges stated, and a count at a single scale cannot decide a threshold statement at all, because the same tile tiles the target at every larger scale.
\textbf{Expectation.} We expect the true set to be a union of schemas of this type, since a threshold or a semigroup appears in all constructions known; whether the thresholds $g_{\min}(e,f)$ and $M^*(e,f)$ have closed forms for all rows is open (for $F_1$, $W$ now has one).
If a core outside all families is ever found, C needs a further family and the classification becomes an open-ended census.
\textbf{Falsifiability.} The cheapest tests of C: the smallest value of $F_5$ outside the known range ($1860=465\cdot2^2$, $4185$, $7440$), the squarefree cores $465$, $966$, $2485$ ($g=1$; C predicts that they are not tile counts), and the first open value of $F_6$ at $g=v$, namely $1617=(4,7,g=7)$.

%=====================================================================
\section{The census}\label{sec:census}

\subsection{What it counts}
The universe is the $392\,074$ non-squarefree integers $N\in[2,10^6]$; $N$ is \emph{excluded} if the decidable hypothesis of the single Lean theorem of the earlier paper fails for it (a program evaluates the hypothesis for each $N$; the theorem is Lean, the evaluation is computation), \emph{known} if some construction gives a tiling, and \emph{undecided} otherwise.
The three sets are disjoint (asserted by the script: no known value is excluded).
The baseline is the census of the previous release.
\begin{center}\small
\begin{tabular}{@{}lrrr@{}}
\toprule
 & excluded & known (tile counts) & undecided\\
\midrule
version 2.2 / release v1.16 & $199\,140$ & $102\,613$ & $90\,321$\\
version 3.0 / release v3.0 & $199\,140$ & $129\,108$ & $63\,826$\\
change & $0$ & $+26\,495$ & $-26\,495$\\
\bottomrule
\end{tabular}
\end{center}
The exclusion side did not change: this release is existence work. The $26\,495$ new known values, by label:
\begin{center}\small
\begin{tabular}{@{}p{4.6cm}rp{6.0cm}@{}}
\toprule
label & values & source\\
\midrule
\PL{} (formal) & $26\,312$ & Theorem~\ref{thm:W} family ($26\,138$ new), collar and tails, macro counts, 13 W-below closures, $k^2$-multiples\\
\CE{} only (exact tiling, no Lean) & $145$ & witness-only values from template searches\\
\PH{} (hand) & $35$ & collar step with $h>1$ on the rows $(1,f),(2,3),(2,5)$ (Theorem~\ref{thm:collar})\\
paper-formal closure over exact seeds & $3$ & upward closure statement of the earlier paper over verified seeds\\
\bottomrule
\end{tabular}
\end{center}
(The $102\,613$ known values of the baseline keep the labels of the previous release.)
Theorem~W alone moves $26\,138$ values from undecided to known: the family has $44\,495$ values $\le10^6$ ($28\,946$ from $W$, $18\,431$ from $T$, union $44\,495$; the largest $f$ is $251$); $18\,357$ of them were already known by other means, none lies in the excluded set.
\textbf{Paliy-claim-only values.} A further $134$ values $\le10^6$ lie in families that Paliy claims (\CP{}; the first are $1380$, $2024$, $3105$, $4554$, $5520$, $8096$) and that we have neither proved nor verified; they are \emph{not} counted as known, they are among the $63\,826$ undecided.

\subsection{Squarefree values}
Of the squarefree $N\le10^6$, $770$ remain open at the exact level ($771$ in version~2.2: the difference is $4830$, a tile count by Paliy's certificate, exact-verified here); at the Lean level the count of the previous release ($791$) is unchanged, since all new Lean theorems concern non-squarefree values.
The open squarefree values $\le1000$ are $231,\ 329,\ 354,\ 385,\ 455,\ 465,\ 517,\ 777,\ 806,\ 935,\ 966$.
The undecided non-squarefree values in $[201,1000]$ are now $38$: $236$, $270$, $280$, $284$, $297$, $332$, $351$, $412$, $420$, $423$, $444$, $480$, $524$, $531$, $552$, $594$, $616$, $644$, $668$, $716$, $732$, $747$, $750$, $816$, $824$, $840$, $855$, $876$, $880$, $908$, $912$, $924$, $927$, $944$, $945$, $956$, $963$, $999$ ($414$ and $639$, tile counts of $W_3(2,5)$ and $T_3(2,5)$, are no longer on the list).
Every $N\le200$ was already decided.

\subsection{What the census is not}
It is not a classification: the undecided values may be tile counts or not. It rests on the families known to us; a tile count in a family we do not know is invisible to it.
The excluded values are excluded by Lean theorems, but the assignment ``this $N$ satisfies the hypothesis'' is made by a program, not by the kernel (as in the earlier paper).
The known values are of the stated labels; only the \PL{} ones are Lean theorems about \emph{that} $N$.

%=====================================================================
\section{Computer evidence on one open family: the $W^{3+}$ census}\label{sec:w3}
This section records exact computations that are \emph{not} proofs and \emph{not} formal.
\subsection{The family and the question}
Let $u<v<2u$ be coprime, $b=v^2-u^2$, tile $(u^2,b,uv)$ (this is the family $F_6$, with $\gamma=2\alpha$), and $\Omega_g$ the isosceles target with base angles $\alpha$, legs $gub$, base $gvb$; a tiling has $g^2b$ tiles.
The earlier papers prove that $\Omega_g$ is not tiled when $(g-1)(v-u)<u$ (Theorem~F) and conjecture that it is tiled only if $v\mid g$ ($V$) and $g\ge g_D(u,v)$ ($A'$); the only tilings known are three Herdt-type counts, $1125$, $19\,600$, $130\,977$.
The first open instances with $v\nmid g$ and $(g-1)(v-u)\ge u$ are $g=3$ for $(u,v)=(4,7)$ (tile $(16,33,28)$, $N=297$) and $(5,8)$ (tile $(25,39,40)$, $N=351$).

\subsection{The words}
In a tiling of $\Omega_g$, the tiles with an edge on a leg partition the leg into segments; reading from the corner $\alpha$ to the apex gives a \emph{word}: the sequence of edge letters ($a$, $b$ or $c$) with the angles of the tiles at their ends.
The length of the leg and the angle sums at the ends leave a finite list of admissible words; hand-proved filters (called $D_2,D_1,K_2,A',G_1$ in the working notes) and a rule for the junction at the apex cut it down to the set we call \emph{$W^{3+}$}.
For each word an exhaustive exact search (integer arithmetic, no floating point) asks whether the part of a tiling near that leg exists.
The verdict ``no tiling'' for a word is therefore conditional on the filters; it covers every tiling whose leg word is that word, and \emph{only} such tilings (it is not a verdict on $\Omega_g$ unless it is returned for every word, including the words outside $W^{3+}$ that the filters exclude).
\begin{center}\small
\begin{tabular}{@{}lrrrr@{}}
\toprule
$(u,v)$, $g=3$ & words & no tiling & tiling found & open\\
\midrule
$(4,7)$ & $191\,940$ (complete list) & $191\,940$ & $0$ & $0$\\
$(5,8)$ & $84\,462$ (complete list) & $84\,462$ & $0$ & $0$\\
$(7,12)$ & $8\,431\,936\,799$ & $8\,014\,982\,151$ ($95.06\%$) & $0$ & $416\,954\,648$\\
\bottomrule
\end{tabular}
\end{center}
For $(4,7)$ and $(5,8)$: all but $323$ and $630$ words fall within $2$ s each; the remainder were rerun at $120$ s, $1800$ s, and without a time limit (the two hardest words of $(4,7)$ took $5614$ s and $3269$ s); every word is decided.
For $(7,12)$ the word list has $8.4\times10^9$ entries and was covered by an adaptive prefix tree (a ``no tiling'' for a prefix refutes all its completions: the justification is that no tile edge straddles the end of a word prefix); the residual $4.2\times10^8$ words are open and the run was stopped because the fraction of prefixes that time out grew with depth.
\textbf{Controls.} The search finds the known tilings when it is given their own words (the words of the tilings of the scaled tile at $m=3,4$, for all three pairs: found and verified), and refutes the reversed words and words refuted independently by a slower program.
\textbf{What this shows and does not show.} No tiling was found, so $V$ and $A'$ are not refuted. It does not show that $297$ and $351$ are not tile counts, because the verdict is for the one-leg word set, not for all tilings and all shapes of target; and without a Lean certificate it is not a theorem. It is evidence, and it narrows where a counterexample could be.


%=====================================================================
\section{Corrections}\label{sec:corr}
Statements made in our working notes, and in one case in an earlier release, that this paper corrects. None of the corrections affects a Lean theorem, a prime-case statement or an excluded value of the census.
\begin{enumerate}
\item \textbf{$183$ is not a tile count.} An earlier draft of the conjecture listed $183$ as a sporadic tile count (attributed to Laczkovich). $183$ is excluded, and every $N\le200$ is decided in Lean; the draft misread a statement about $5\cdot183^2$. The Lean predicate $F_4$ never contained it.
\item \textbf{Tiles of $T_M$ and $W_M$ below the scale $f$ exist for $e\ge2$.} Our notes asserted that $T_M(e,f)$ and $W_M(e,f)$ have no tiling for $M<f$ when $e\ge2$ (``$B_1$''), and that the $T$ half of this statement stood after the $W$ half failed. Both are false: Theorem~\ref{thm:W} and the glue of Section~\ref{sub:Tcor}. The version 1 paper \cite{ReleasePaper} recorded, as an observation, that ``in every case decided so far'' the scale $f$ is the threshold; Corollary~\ref{cor:W435} is a counterexample ($W_4(3,5)$, $M=4<5$). The open problems of that paper on the scales $f/(2e)<M<f$ are partly closed here (Section~\ref{sec:thmW}).
\item \textbf{The template ``exists iff $4f<9e$'' is false.} An intermediate note (13 witnesses below $f$, then a replay of one search) stated that the six-piece template exists only for $4f<9e$, was absent for $432$ coprime pairs with $f\le60$, and that the conjectured threshold held only there. This was an artefact: the replay of one search history has an empty cone for the other pairs, but the geometry of the same six pieces works for all $2\le e\le f-2$ (Theorem~\ref{thm:W}; for $(2,5)$, $(3,7)$, $(4,9)$, $(5,12)$, $(3,8)$ the tilings exist).
\item \textbf{The size conjecture is false.} The working notes (and a conjecture in a longer draft) stated that a target not similar to the tile has longest side at least $c^2=f^4$. $W_4(3,5)$ has longest side $500<625$.
\item \textbf{Paliy's $4830$ is a tiling.} Our notes said that his $4830$-tile file is not a tiling. We had checked the wrong file (his refuted moment counterexample, not his certificate). His certificate, $4830$ copies of $(24,11,31)$ on the triangle $(961,1155,1426)$, is verified exactly (Section~\ref{sec:macro}).
\item \textbf{The $\theta$-isosceles statement.} Our notes stated ``targets $\theta_t(e,f)$ with $te\ge f$ have no tiling''. That is false: $\theta_4(1,2)$ ($te=8$) and $\theta_8(2,3)$ are tiled. The proved statement is the necessary bound $te\ge2f$; the true statement is a threshold (a semigroup of scales, Paliy), not an exclusion.
\item \textbf{Priority of the $\alpha$ up-set.} The up-set in $g$ of the base-$\alpha$ rows and the form $b\kappa K^2$ are Paliy's, dated 29 September 2026 (Section~\ref{sec:collar}); our notes had presented the up-set as ours. The scale additivity of the rows is also his.
\item \textbf{Counts in a note.} The census effect of the collar step was first stated as $235$ values; the correct figure, by an audit that recounted, was $194$ ($119$ in the rows $(1,f),(2,5)$ and $75$ in $(2,3)$); the figures of this release supersede both.
\item \textbf{The threshold formula $(M-1)(f-e)\ge f$.} Used as a threshold for all rows it contradicts the proved side theorem $W_M$ none if $Me<f$ when $e/f<(3-\sqrt5)/2$ (for instance $(2,7)$: the formula gives $M=3$, but $3\cdot2<7$). The correct threshold is $M^*(e,f)=\max(\lceil f/e\rceil,\lceil f/(f-e)\rceil+1)$.
\item \textbf{The pattern $g_{\min}=f+1$.} Our notes observed that the first realised scale of a base-$\alpha$ row was $f+1$ for even $f$ and recorded this as a pattern; $g_{\min}(2,5)$ and $g_{\min}(1,6)$ are $\le2$ (Theorem~\ref{thm:tails}).
\end{enumerate}
The register of corrections of the project is in the repository (file \texttt{TIMELINE.md} records the dates).

%=====================================================================
\section{What is not proved}\label{sec:notproved}
\begin{itemize}
\item \textbf{The exclusion half of Conjecture~C.} No description of the set of tile counts above $200$ is proved; the census has $63\,826$ undecided non-squarefree values and $770$ squarefree values $\le10^6$.
\item \textbf{The no side of Theorem~W for $e>f/2$:} for each row exactly the scales $M$ with $(M-1)(f-e)<f<(M+1)(f-e)$ (at most two) are undecided; and \textbf{$e=1$ and $e=f-1$ for $T_M$} (first case $T_3(1,4)$); the $T$ side in general (only $2Me<f$ is excluded).
\item \textbf{$V$ and $A'$ for $\gamma=2\alpha$:} the evidence of Section~\ref{sec:w3} is not a proof; $(7,12)$ is incomplete ($95.06\%$).
\item \textbf{$\Theta$ and $S_3$:} the thresholds of the $\theta$-isosceles family and of the base-$\alpha$ rows; the squarefree cores $465$, $966$, $2485$; the values $1860$, $4185$, $7440$.
\item \textbf{$15m^2$ for every $m\ge4$ as a Lean theorem}; the collar step for arbitrary height $h$ (hand only); Paliy's claims that we have not verified ($134$ values $\le10^6$ in the census; the mixed-corner family beyond $4830$).
\item \textbf{The squarefree values:} $770$ open; the squarefree tile counts known to us are of the classical forms, or come from a short list of sporadic and layout values (the $22$ of the earlier release, $77$, $143$, $4830$).
\item \textbf{Outside check.} The Lean statements and the proofs have been compared and audited only inside this project, including the statement of \texttt{theoremP\_6} (the prime case of the earlier releases) and of Theorem~W; a check by a reader outside the project has not been made, and announcing the prime case as proved remains the author's decision.
\item The search that found the macro layouts is not part of this release and is not described further than in Section~\ref{sec:macro}; its outputs are verified independently.
\end{itemize}

%=====================================================================
\section{Credits, with dates}\label{sec:credits}
``Ours'': statement and proof produced in this work. ``Formalized'': another author's statement and construction, our Lean proof. Dates are of first public record (repository commit dates, arXiv versions) unless marked.
\begin{center}
\small
\begin{tabular}{@{}p{4.4cm}p{1.8cm}p{6.4cm}@{}}
\toprule
result & status here & whose, and when\\
\midrule
Prime counts; the base-$\beta$ reduction & earlier releases & Paliy, first public proof 2026-09-29 \cite{PaliyPrime}; ours independent; exclusion of $e=1$ public 2026-09-12 (Zenodo v3.2); Beeson's withdrawal note for arXiv:2607.19572 (v2, 2026-09-24) \cite{BeesonWithdrawn} credits the author, cited here as public record\\
Targets $W_M$, $T_M$; tilings for $M\ge f$ & formalized & Paliy ($W/\beta$ framework, caps, collars: 2026-10-03; all $M\ge v$: 2026-10-10); Beeson and Ro\c{s}u (the targets, the seed)\\
\textbf{Theorem~W}, six-piece template & ours & template found by our search; the setting and the earlier thresholds are Paliy's and Beeson's; the sharpness side ($Me<f$) is from the earlier releases\\
$\alpha$ up-set, form $b\kappa K^2$, scale additivity & Paliy's; Lean here & Paliy, 2026-09-29 \cite{PaliyRepo}; seeds $84$, $189$: Beeson (Figs.~28--29), $189$: Harries\\
Collar step, parametric seed $e=1$, tails & ours & independent of Paliy's construction; credit for the up-set as above\\
Macro pieces & ours (layouts, Lean) & macro-dissections: Laczkovich; Herdt ($1125$); seeds $15m^2$: Paliy (2026-10-08)\\
$15m^2$ for all $m\ge4$, $4830$, mixed-corner family & Paliy's (\CP{}); $4830$ verified here & Paliy, 2026-10-08\\
$990$, $264$, $897$, $240$, $312$ & earlier releases & Paliy: $990$ on 2026-10-06, $264$ on 2026-10-07; the others as dated in the earlier releases\\
$\theta$-isosceles family, scale additivity & Paliy's (\CP{}) & Paliy, 2026-09-29\\
Run induction, Theorem~F and $D'$ & earlier releases & method: Beeson and Ro\c{s}u (preprint of 2026-10-16); $g=1$: Beeson 2012\\
$1125$-tiling & earlier releases & Herdt, published by Beeson\\
$88$ & earlier releases & Harries\\
$2\pi/3$ families $F_3$, $15m^2,33m^2$ & formalized & Zhang (arXiv:2512.22696); seeds Paliy, Herdt\\
Rep-tiling and $F_0$ & --- & Snover, Waiveris and Williams\\
Classification of pairs (triangle, tile) & --- & Laczkovich\\
Conjecture~C, census, $W^{3+}$ evidence & ours & ---\\
\bottomrule
\end{tabular}
\end{center}
The statements ``Paliy claims'' in this paper are statements in his public repository \cite{PaliyRepo}; we have read his documents and checked some examples, and we say so where we did. Nothing here is a claim of priority over him for the results listed as his; the dates are given so that the order is visible.


%=====================================================================
\appendix
\section{Proof of Theorem~W}\label{app:W}

Credit: the targets $W_M,T_M$ and the framework are Beeson's and Paliy's; the six-piece template was found by our search, and every identity below was also confirmed by a computer algebra system and by exact integer computation; Lean checks the whole argument (Section~\ref{sec:thmW}).
Throughout, $e,f$ are coprime, $2\le e\le f-2$, $g=f-e\ge2$, and $a,b,c,\kappa$ are as in Section~\ref{sub:tile}.

\subsection{Complex coordinates}
Put $w=\sqrt{4f^2-e^2}$ and write the complex number $x+iwy$ as $[x,y]$. Then $\mathrm{Area}(\tile)=ebw/4$, and
\[
\varepsilon=e^{i\alpha}=\frac{\kappa+iew}{2f^2},\qquad \eta=e^{i\beta}=\frac{e(3f^2-e^2)+ibw}{2f^3},\qquad \omega=\overline{\varepsilon\eta}=e^{-i(\alpha+\beta)} .
\]
Polynomial identities (clear denominators and expand): $\varepsilon^3\eta^2=-1$ (this is $3\alpha+2\beta=\pi$); $a+b\omega=c\bar\eta$ (the tile closes up); $f\omega+e\varepsilon=f\bar\eta$.
A triangle similar to $\tile$ with sides $sa,sb,sc$ is a union of $s^2$ tiles, and a parallelogram with sides $sa$, $uc$ and angle $\beta$ is a union of $2su$ tiles (Proposition~\ref{prop:ladder}).

\subsection{The template}
Let $M,t,r,q$ be positive integers and put
\[
v=fM-er,\quad s_1=f(r-t),\quad s_2=fr-eM,\quad x=ft-eq,\quad u_1=fq-et,\quad u_2=fq-er .
\]
Points (vector sums in $\mathbb C$):
\begin{align*}
O&=0, & B&=fbM, & A&=f^3M\varepsilon^2, & P&=fbM\varepsilon,\\
P_2&=b(fM-et)\varepsilon, & P_3&=P-eat\omega, & Q&=P+fbt\omega, & R&=P+fbr\omega,\\
C'&=bv\varepsilon, & B'&=cv, & V_2&=P+f^3t\bar\eta-ef^2r, & X_4&=P-fbq\varepsilon,\\
U_1&=P-ef^2t+ef^2q\bar\eta, & U_2&=P-ef^2r+ef^2q\bar\eta .
\end{align*}
The target is the triangle $OAB$: $|OB|=Mfb$, $|OA|=Mf^3$, $|AB|=Me\kappa$, with angles $2\alpha$ at $O$, $\beta$ at $A$, $\alpha+\beta$ at $B$; that is $W_M(e,f)$. The point $P$ lies on $AB$ ($A-P=-afM\omega$, $B-P=ebM\omega$) and $OP=OB$.
The pieces, counter-clockwise:
\begin{center}\small
\begin{tabular}{@{}lllp{5.3cm}l@{}}
\toprule
 & vertices & kind & description & tiles\\
\midrule
$\pi_0$ & $A,O,P_2,P_3$ & trapezoid & triangle $O,A,P$ (scale $fM$) minus the corner triangle $P,P_2,P_3$ (scale $et$) & $f^2M^2-e^2t^2$\\
$\pi_1$ & $O,B',C'$ & triangle & scale $v$ & $v^2$\\
$\pi_2$ & $P_3,P_2,Q$ & triangle & scale $ft$ & $f^2t^2$\\
$\pi_3$ & $B,Q,V_2,B'$ & trapezoid & triangle $R,Q,V_2$ (scale $s_1$) minus the corner $R,B,B'$ (scale $s_2$) & $s_1^2-s_2^2$\\
$\pi_4$ & $Q,U_1,U_2,V_2$ & parallelogram & sides $a s_1$ ($V_2\to Q$) and $cx$ ($Q\to U_1$) & $2s_1x$\\
$\pi_5$ & $U_1,P_2,C',U_2$ & trapezoid & triangle $X_4,U_1,P_2$ (scale $u_1$) minus the corner $X_4,C',U_2$ (scale $u_2$) & $u_1^2-u_2^2$\\
\bottomrule
\end{tabular}
\end{center}
The area identity $\sum(\text{tiles})=M^2(2f^2-e^2)$ is an identity in $\mathbb Z[e,f,M,t,r,q]$ (checked by computer algebra).

\subsection{Lemma A (shapes)}
\emph{If $s_1>s_2>0$, $u_1>u_2>0$, $v,ft>0$ and $fM>et$, each piece is the stated union of tiles.}
\begin{itemize}
\item $\pi_0$: $OP=bfM\varepsilon$ and $OA=cfM\varepsilon^2$ have lengths $bfM$, $cfM$ enclosing $\alpha$: SAS gives a triangle of scale $fM$. $P_2$ lies on $OP$ and $P_3$ on $PA$ with $PP_2:PO=PP_3:PA=et:fM$, so $P_2P_3\parallel OA$ and the cut-off triangle $PP_2P_3$ has scale $et$.
\item $\pi_1$: $OB'=cv$ and $OC'=bv\varepsilon$ enclose $\alpha$: scale $v$.
\item $\pi_2$: $P_3Q=(ea+fb)t\omega=f^3t\omega$ has length $c\cdot ft$, and $|P_3P_2|=c\cdot et=a\cdot ft$; the angle at $P_3$ is the angle $\beta$ of the triangle $PP_2P_3$, similar to $OPA$: scale $ft$.
\item $\pi_3$: $Q-V_2=as_1$ is real (this is the closing identity $a+b\omega=c\bar\eta$ times $ft$); $R-Q=-bs_1\omega$; the angle at $Q$ between the directions $\omega$ and $\pi$ is $\gamma$: sides $a,b$ enclosing $\gamma$, scale $s_1$. $B=R-bs_2\omega$ lies on $RQ$ and $B'$ lies on $RV_2$ with $RB':RV_2=s_2:s_1$.
\item $\pi_4$: $U_1-Q=-f^2x\bar\eta=U_2-V_2$, so $V_2QU_1U_2$ is a parallelogram with sides $as_1$, $cx$ and angle $\beta$ at $Q$.
\item $\pi_5$: $P_2-X_4=bu_1\varepsilon$; $|U_1-X_4|=cu_1$, $|U_1-P_2|=au_1$; $C'$ lies on $X_4P_2$ with $X_4C':X_4P_2=u_2:u_1$, and $U_2$ on $X_4U_1$ with the same ratio.
\end{itemize}
Each statement is a polynomial identity in $e,f,M,t,r,q$ (clear denominators) and was checked by computer algebra; the Lean proof has one lemma for each. $\square$

\subsection{Lemma B (geometry reduces to inequalities)}
All six pieces are convex and counter-clockwise when the scales are positive (their doubled signed areas are the tile counts times a positive constant). Hence a piece lies in the target iff its vertices do, and two pieces have disjoint interiors if one edge line of one has the other piece on its closed outer side; the edges needed are: $\pi_0|\pi_1$: $OP_2$; $\pi_0|\pi_2$: $P_2P_3$; $\pi_0|\{\pi_3,\pi_4,\pi_5\}$: $OP_2$; $\pi_1|*$: $B'C'$; $\pi_2|*$: $P_2Q$; $\pi_3|\{\pi_4,\pi_5\}$: $QV_2$; $\pi_4|\pi_5$: $U_1U_2$.
With the area identity this gives a tiling (closed convex pieces inside the closed convex target of total area equal to its area cover it).
Every cross product needed is a product of positive constants (depending on $e,f$) and factors from the list below; the sign of the constant prefactors is read once, at one instance, and is independent of the parameters:
\begin{itemize}\itemsep0pt
\item C2: $r-t>0$;\quad C3: $eM-ft>0$;\quad C4: $fr-eM>0$;\quad C5: $fM-er>0$;
\item C6: $fq-er>0$;\quad C7: $ft-eq>0$;\quad C8: $f^2t-e^2r>0$;\quad also $M>q$, $t>0$, $r>0$;
\item LA: $bfM-e\kappa t+f^3q>0$;\quad LB: $bfM-e\kappa r+f^3q>0$;
\item LC: $efbM-e^2\kappa r+f^4t>0$
\end{itemize}
(together with three expressions that are positive for any positive parameters: $bM+eft$, $eMf+bt$, $br+e^2t$; LB implies LA since $r>t$).
\emph{Lean:} one lemma per factor (modules \texttt{ThmWIneq0}--\texttt{ThmWIneq4}).

\subsection{Lemma C (the parameters satisfy the list)}
Assume (H1), (H2). Put $n=\lfloor eM/f\rfloor$, $\rho=eM-fn\in[0,f)$, $s=M-n$, $\delta=1$ if $\rho=0$ and $\delta=0$ otherwise, and
\[
t=n-\delta,\qquad r=n+1,\qquad q=n .
\]
\textbf{Facts.}
(F1) $gM=fs-\rho$, and (H2) gives $gM\ge f+g$; so $gM/f>1$ and $s=\lceil gM/f\rceil\ge2$.
(F2) $n\ge1$ by (H1).
(F3) If $\rho=0$ then $f\mid M$ (coprimality), $M=fm'$, $n=em'\ge e\ge2$.
(F4) $eg\ge f+1$: $eg-f=(e-1)(g-1)-1\ge1$, since $e,g\ge2$ and not both equal to $2$ (coprime).
(F5) $\boldsymbol{gn\ge e+1}$. If $g\ge e+1$: $gn\ge g\ge e+1$. If $g\le e$: write $f=gh+j$ with $0<j<g$ (coprimality of $g$ and $f$) and $h\ge2$; then $e=g(h-1)+j$, and (H2) gives $M-1\ge\lceil f/g\rceil=h+1$; now $e(h+2)-hf=g(h-2)+2j>0$, so $eM>hf$, $n\ge h$, and $gn\ge gh=f-j\ge f-g+1=e+1$.
\textbf{The factors.}
\begin{itemize}
\item $t\ge1$: for $\delta=0$, $t=n\ge1$; for $\delta=1$, $t=n-1\ge e-1\ge1$ by (F3). And $r>t$, $r>0$, $q=n>0$. $M>q$ because $n\le eM/f<M$. (C2 holds.)
\item C3: $ft\le fn\le eM$, with equality only if $\delta=0$ and $\rho=0$, which is impossible; for $\delta=1$, $ft=fn-f<eM$.
\item C4: $fr-eM=f-\rho>0$.
\item C5: $fM-e(n+1)\ge(bM-ef)/f>0$ because $bM=(f+e)gM\ge(f+e)(2f-e)=2f^2+ef-e^2>ef$ (by F1: $gM\ge f+g=2f-e$).
\item C6: $fq-er=gn-e\ge1$ by (F5).
\item C7: for $\delta=0$, $ft-eq=gn>0$; for $\delta=1$, $ft-eq=f(n-1)-en=gn-f\ge ge-f\ge1$ by (F3), (F4).
\item C8: for $\delta=0$, $f^2t-e^2r=bn-e^2=(f+e)gn-e^2\ge(f+e)(e+1)-e^2>0$; for $\delta=1$, it is $bn-f^2-e^2\ge be-f^2-e^2=eg(f+e)-f^2-e^2\ge(f+1)(f+e)-f^2-e^2>0$.
\item LB: substituting $M=(fs-\rho)/g$ and $n=(eM-\rho)/f$ one finds the polynomial identity $\mathrm{LB}=s\Psi-\kappa(\rho+e)$ with $\Psi=f^3+2ef^2-e^2f-e^3>0$ (checked by computer algebra). With $s\ge2$ and $\rho\le f-1$: $\mathrm{LB}\ge2\Psi-\kappa(f+e)+\kappa>2\Psi-\kappa(f+e)=e(2f+e)(f-e)>0$. Then LA$=\mathrm{LB}+e\kappa(r-t)>0$.
\item LC: $\mathrm{LC}=efbM+b^2n-e^2\kappa-\delta f^4$ (using $f^4-e^2\kappa=b^2$). For $\delta=0$: $efbM=ef(f+e)gM\ge ef(f+e)(2f-e)$, so $\mathrm{LC}\ge e\,[f^2(2f-e)-e^2(f-e)]>0$. For $\delta=1$: $M=fm'$, $n=em'$, $\mathrm{LC}=m'eb\kappa-e^2\kappa-f^4\ge e\kappa(b-e)-f^4$, and $e(b-e)=eg(f+e)-e^2\ge2(f-2)(f+2)-(f-2)^2=f^2+4f-12>f^2$ (using $eg\ge2(f-2)$ for $2\le e\le f-2$, $f+e\ge f+2$, $e\le f-2$), while $\kappa>f^2$; so $\mathrm{LC}>0$.
\end{itemize}
Positivity of the scales: $v=fM-er>0$ is C5, $s_1=f(r-t)>0$ is C2 times $f$, $s_2=fr-eM>0$ is C4, $s_1>s_2$ is $f(r-t)>fr-eM$, i.e.\ C3, $x=ft-eq>0$ is C7, $u_1=fq-et>0$ follows from C6, $u_1>u_2$ is $e$ times C2, $u_2=fq-er>0$ is C6, and $fM>et$ follows from C3.
Lemmas A, B and C prove Theorem~\ref{thm:W} for $W_M$; the count of tiles is $M^2(2f^2-e^2)$ by the area identity. The inequality $4f<9e$ is not used anywhere. $\square$

\subsection{Evidence that the proof is right}
See Section~\ref{sub:Wevidence}. In addition, the parameter triple $(t,r,q)$ is not unique (for instance $q=n+3$ also works at $(7,13,13)$), and the canonical triple satisfies the whole list exactly when $M\ge M^*(e,f)$, for all coprime pairs with $f\le40$ and $M<4f$: the bound (H1)+(H2) is exactly the range in which the canonical parameters work.

%=====================================================================
\section{Data: the explicit tilings}\label{app:data}
\subsection{The thirteen first witnesses with $M<f$}
Found first by a computer search, and now instances of Theorem~\ref{thm:W}; every one is a Lean theorem (module \texttt{Draft/DebtBWTilable}, with the closure upward in $M$).
\begin{center}\small
\begin{tabular}{@{}lrr@{\qquad}lrr@{}}
\toprule
$(e,f,M)$ & $W_M$ tiles & $T_M$ tiles & $(e,f,M)$ & $W_M$ tiles & $T_M$ tiles\\
\midrule
$(3,5,4)$ & $656$ & $1056$ & $(7,12,4)$ & $3824$ & $6128$\\
$(4,7,4)$ & $1312$ & $2096$ & $(7,10,5)$ & $3775$ & $6275$\\
$(5,8,4)$ & $1648$ & $2672$ & $(7,9,6)$ & $4068$ & $6984$\\
$(5,7,5)$ & $1825$ & $3050$ & $(8,11,5)$ & $4450$ & $7475$\\
$(5,11,3)$ & $1953$ & $3042$ & $(9,11,7)$ & $7889$ & $13\,818$\\
$(5,9,4)$ & $2192$ & $3488$ & $(6,11,4)$ & $3296$ & $5232$\\
$(7,11,4)$ & $3088$ & $5024$ & & &\\
\bottomrule
\end{tabular}
\end{center}
\subsection{The shipped files}
Directory \lean{data/witnesses/v3_exact}: $97$ files, each named by the number of tiles, the kind and the parameters (for example \lean{tiling_414_W_2_5_M3.exact.txt.gz}). Each begins with a comment line stating the number of tiles and the side lengths of the tile; every other line gives the three vertices of one tile in exact rational coordinates $(X,Y\sqrt D)$. Kinds: \texttt{W}, \texttt{T} (Theorem~W and its glue: $37$ and $15$ files), \texttt{collar} ($6$) and \texttt{macro} ($38$) (Sections~\ref{sec:collar}, \ref{sec:macro}), \texttt{F3\_paliy} (Paliy's $4830$). Verification: the exact verifier of the repository (\texttt{bte verify}) reports, for each file, that the tiles are congruent to the stated tile, that their interiors are pairwise disjoint, that the total area equals the area of the hull, and that the hull is a triangle. All $97$ pass.

\section{Index of the Lean statements}\label{app:lean}
All in the library of the repository (directory \texttt{lean/}); aggregate \lean{Erdos634.AllV3} (the modules new in this release; build one module at a time through \texttt{code/lakeserial.py}); the axioms of $98$ of them are printed by \texttt{lean/axioms/AxiomsAllV3.lean} and every line reads \texttt{[propext, Classical.choice, Quot.sound]}.
\begin{center}\small
\begin{tabular}{@{}p{4.9cm}p{8.1cm}@{}}
\toprule
statement & Lean name (namespace \texttt{Erdos634.}\dots)\\
\midrule
Theorem~\ref{thm:W} & \lean{Draft.ThmW.thmW}; counts \lean{Draft.ThmW.tileCount_W}, \lean{tileCount_T}; examples \lean{inst_414}, \lean{inst_656}, \lean{inst_639}, \lean{inst_27324}\\
13 first witnesses, closures & \lean{Draft.DebtB.closure_656} etc.; \lean{DebtB.counts_of_wTilable}\\
Collar layer step & \lean{Draft.CollarLayer.layer_core}, \lean{Draft.CollarFam.layer_step}\\
Theorem~\ref{thm:seedA} & \lean{Draft.CollarUp.tileCount_e1_up}; instances \lean{tileCount_189}, \lean{tileCount_16740}, \lean{tileCount_159040}\\
Theorem~\ref{thm:tails} & \lean{Draft.CollarFam.tileCount_465_sq}, \lean{Draft.CollarFam5.tileCount_2485_sq5}, \lean{Draft.CollarFam17.tileCount_966_sq17}\\
Macro pieces (Thm~\ref{thm:macro}) & \lean{Draft.MacroPiece.macroTFam_gen}; \lean{Draft.TM375.tileCount_375} etc.; \lean{Draft.T62125.tileCount_62125}; \lean{Draft.AllMacro.tileCount_78246}\\
Twelve square multiples & \lean{Draft.TemplateSquares.tileCount_46500} etc.\\
Conjecture C, existence half & \lean{ConjC.Exist.conjC_existence}; \lean{ConjC.BetaCount.tileCount_beta_ge}\\
\bottomrule
\end{tabular}
\end{center}

\section{Reproducing the checks}\label{app:repro}
From a clean clone, with the toolchain of \texttt{lean/lean-toolchain}: build the aggregate one module at a time, print the axioms, and verify the exact tilings:
\begin{verbatim}
python3 code/lakeserial.py lean build Erdos634.AllV3
(cd lean && lake env lean axioms/AxiomsAllV3.lean) \
   | python3 code/axiom_gate.py
BTE=rust/bonfioli_tiling_engine/target/release/bte
for f in data/witnesses/v3_exact/*.gz; do
  $BTE verify "$f" | grep -q VERIFIED || echo FAIL "$f"
done
\end{verbatim}
The build needs about $6$ GB of memory for its largest module and about $22$ GB of disk for the build products; one module is built at a time. The scripts \texttt{code/reproduce.sh} (all earlier checks) and the above give the complete set.

%=====================================================================
\begin{thebibliography}{99}
\bibitem{BeesonI} M.~Beeson, \emph{Triangle tiling I: The tile is similar to ABC or has a right angle}, arXiv:1206.2231.
\bibitem{BeesonRosu} M.~Beeson and G.~Ro\c{s}u, \emph{Tiling a triangle into a prime number of congruent triangles}, preprint, version of 16 October 2026.
\bibitem{BeesonIII} M.~Beeson, \emph{Triangle tiling III: the case $3\alpha+2\beta=\pi$}, arXiv:1206.2229.
\bibitem{BeesonIso} M.~Beeson, \emph{Tilings of an isosceles triangle}, arXiv:1206.1974.
\bibitem{BeesonWithdrawn} M.~Beeson, \emph{No prime tiling of an isosceles triangle}, arXiv:2607.19572 (withdrawn in version 2, 24 September 2026, with a note attributing the theorem to V.~Bonfioli).
\bibitem{Erdos633} M.~Beeson, M.~Laczkovich and Y.~X.~Zhang, \emph{Solution of Erd\H{o}s Problem~633}, arXiv:2604.03609.
\bibitem{Laczkovich1995} M.~Laczkovich, \emph{Tilings of triangles}, Discrete Math.\ \textbf{140} (1995), 79--94.
\bibitem{Laczkovich2012} M.~Laczkovich, \emph{Tilings of convex polygons with congruent triangles}, Discrete Comput.\ Geom.\ \textbf{48} (2012), 330--372.
\bibitem{PaliyPrime} D.~Paliy, \emph{No scale-one triangle in either rational Group-1 family}, and the deduction of the prime case, repository \texttt{github.com/DenisUkranian/erdos-634-tilings}, files \texttt{docs/prime-case-candidate.md} and \texttt{docs/prime-case-dependencies.md}, first published 29 September 2026; unrefereed.
\bibitem{PaliyRepo} D.~Paliy, repository \texttt{github.com/DenisUkranian/erdos-634-tilings}: notes on eventual rational families (\texttt{docs/eventual-rational-families.md}), the $\theta$-branch (\texttt{docs/theta-branch.md}), the $W/\beta$ framework and the macro-dissection files, September--October 2026; unrefereed.
\bibitem{ReleasePaper} V.~Bonfioli, \emph{Congruent dissections of a triangle: prime counts, the counts up to $200$, and a family of isosceles targets}, release v1.13.1, Zenodo, October 2026, doi:10.5281/zenodo.23400433.
\bibitem{Release22} V.~Bonfioli, \emph{Congruent dissections of a triangle: runs, windows and the corner}, version 2.2, release v1.16, Zenodo (all versions), doi:10.5281/zenodo.23400433.
\bibitem{SWW} S.~L.~Snover, C.~Waiveris and J.~K.~Williams, \emph{Rep-tiling for triangles}, Discrete Math.\ \textbf{91} (1991), 193--200.
\bibitem{Soifer} A.~Soifer, \emph{How does one cut a triangle?}, 2nd ed., Springer, 2009.
\bibitem{Zhang} Y.~X.~Zhang, \emph{Tiling triangles with $2\pi/3$ angles}, arXiv:2512.22696.
\end{thebibliography}

\end{document}
